\documentclass[10pt,a4paper]{amsart}
\usepackage[margin=19mm]{geometry}
\usepackage{amsmath,amssymb,mathtools}
\usepackage[scr=rsfs,cal=euler]{mathalfa}
\usepackage{xfrac}
\usepackage[unicode,hidelinks,bookmarks=false]{hyperref}
\newcommand{\ie}{i.e.\ }
\newcommand{\cf}{cf.\ }
\newcommand{\resp}{respectively\ }
\newcommand{\Subsection}{\subsection}
\providecommand{\nobreakdash}{\nobreak\hbox{-}}
\setlength{\emergencystretch}{2em}
\multlinegap0pt
\usepackage{amssymb}
\newtheorem{lemma}{Lemma}
\def\R{{\mathbb R}}
\newcommand*{\qefed}{=\mathrel{\vcenter{\baselineskip0.5ex \lineskiplimit0pt

                     \hbox{\scriptsize.}\hbox{\scriptsize.}}}%
                     }
\title{Resolvent estimates for the Schr\"odinger operator with $L^\infty$ electric and magnetic potentials and applications to the local energy decay}
\author{Andr\'es Larra\'in-Hubach, Jacob Shapiro,, Georgi Vodev}
\date{Source: arXiv:2604.09936v1 (CC BY 4.0)}
\begin{document}
\maketitle

\noindent\emph{Source and licence.} Resolvent estimates for the Schr\"odinger operator with $L^\infty$ electric and magnetic potentials and applications to the local energy decay. Version arXiv:2604.09936v1 (CC BY 4.0), distributed by arXiv under the Creative Commons Attribution 4.0 International licence, http://creativecommons.org/licenses/by/4.0/, as recorded in the arXiv licence metadata for this exact version and shown on the abstract page https://arxiv.org/abs/2604.09936v1. Full licence notice: \texttt{licenses/arxiv-2604.09936-CC-BY-4.0.txt}.\par\smallskip

\noindent\textit{Editorial scope.} Counted unit: the article's lemma 2.2 (source0.tex lines 431--438). The article's complete original proof of it is delivered with it (source0.tex lines 506--539, one \texttt{proof} environment of 34 lines, which the article places away from the statement). No mathematics is written, repaired or re-derived: the statement and the proof are the article's own lines, in the article's own notation, and no retained line is substituted or re-flowed.  Retained uncounted as the source-local context the proof needs: lines 162--164.  Nothing was omitted from any retained span, so the package carries no omission record.  No retained line contains a citation, so the wrapper carries no bibliography and no bibliography entry.  No retained line contains a figure environment, an image-inclusion command or drawing code, so no image is delivered and the package declares no asset dependency.  Editorial formatting only, in addition: an AMS \texttt{amsart} preamble that restates the article's own declarations and macros used by the retained lines, namely the environments \texttt{lemma} on separate counters, together with the standard packages those lines need.  The retained environments print in this document as Lemma 1 in order of appearance, whereas the article prints the same items with the numbering of its own full document.  The complete native source of the article as distributed in the arXiv e-print for this exact version is bundled unchanged as \texttt{sources/198235/source0.tex}.
\begin{equation}\label{eq:1.7}
\left|\partial_r^j\Theta(r)\right|\le\widetilde C\Theta(r),\quad\forall r\ge 0,\quad j=1,2,
\end{equation}

 \begin{lemma} \label{2.2} 
Let $\alpha$ be a multi-index such that $|\alpha|\le 2$. 
 Then, there exist constants $C > 0$ and $\theta_0 > 0$ such that for all $\theta \ge \theta_0$ we have the estimate
\begin{equation}\label{eq:2.11}
\left\|\mu\partial_x^\alpha\left(P_0-i\theta^2\right)^{-1}\mu^{-1}
\right\|_{L^2\to L^2}\le C \theta^{|\alpha|-2}.
\end{equation}
\end{lemma}

 \begin{proof}[Proof of Lemma \ref{2.2}] 
We utilize the standard resolvent estimate: there exists $C > 0$ so that for any multi-indices $\alpha, \beta$ with $|\alpha| + |\beta| \le 2$, and for any $\theta \ge 1$, 
\begin{equation} \label{free resolv bd}
\| \partial_x^\alpha (P_0 - i\theta^2)^{-1} \partial_x^\beta \|_{L^2 \to L^2} \le C \theta^{|\alpha| + |\beta| - 2}.
\end{equation}
 We compute the commutator
\begin{equation*}
\begin{split}
\mu(P_0 - i \theta^2) - (P_0 - i\theta^2) \mu &= - \Delta \mu -2 (\nabla \mu) \cdot \nabla  \\
&=- \sum_{j =1} \partial_{x_j} (\partial_{x_j} \mu) - \Delta \mu \qefed [ \mu, P_0],
\end{split}
\end{equation*}
which implies the resolvent identity
\begin{equation*}
\mu (P_0 - i\theta^2)^{-1} = (P_0 - i \theta^2)^{-1}  \mu -  (P_0 - i \theta^2)^{-1}  [ \mu, P_0]  (P_0 - i \theta^2)^{-1} .
\end{equation*}
Therefore, for $f \in C_0^\infty(\R^d)$, 
\begin{equation} \label{resolv id with f}
\mu(P_0 - i \theta^2)^{-1} \mu^{-1} f = (P_0 - \theta^2)^{-1} f -  (P_0 - \theta^2)^{-1} [\mu, P_0] (P_0 - i\theta^2)^{-1} \mu^{-1} f.
\end{equation}
Continuing, for any $u \in C_0^\infty(\R^d)$, 
\begin{equation} \label{commutator insert weights}
[\mu, P_0] u = \left[ -(\Delta \mu) \mu^{-1} - \sum_{j=1}^d \partial_{x_j} \big((\partial_{x_j} \mu) \mu^{-1} \big)  \right] (\mu u)
\end{equation}
By \eqref{eq:1.7}, $(\partial_{x_j} \mu)\mu^{-1}$ and $(\Delta \mu) \mu^{-1}$ are bounded multiplication operators on $L^2$. So \eqref{commutator insert weights} continues to hold for any $u \in L^2(\R^d)$, with both sides being interpreted as members of the negative index Sobolev space $H^{-1}(\R^d)$. Therefore, by \eqref{free resolv bd}, \eqref{resolv id with f}, and \eqref{commutator insert weights}, 

\begin{equation*}
\begin{split}
\| \mu (P_0 - i\theta^2)^{-1} \mu^{-1} f \|_{L^2} &\le \| (P_0 - i\theta^2)^{-1} f \|_{L^2} + \| (P_0 - i\theta^2)^{-1} \|_{H^{-1} \to L^2} \| \mu (P_0 - i\theta^2)^{-1} \mu^{-1} f \|_{L^2} \\
& \le C(\theta^{-2} \| f\|_{L^2} + \theta^{-1} \| \mu (P_0 - i\theta^2)^{-1} \mu^{-1} f \|_{L^2}),
\end{split} 
\end{equation*}
for some $C > 0$ independent of $\theta$ and $f$. Thus taking $\theta$ large enough we can absorb the second term on the right side into the left hand side, completing the proof of \eqref{eq:2.11} when $|\alpha| = 0$. Standard arguments using elliptic regularity then imply the general case  $|\alpha | \le 2$.\\
 \end{proof}

\end{document}
